\scanpage{175}
\mspage{91}
% archive sheet number: 174
% lines 1-7, continuation
auxdites sommes \deleted{\unsure{implique} \unsure{(|$C$}} (chose également \inserted{déjà} acquise, puisque $\operatorname{card}(\operatorname{Hom}(j,j'))\leq\pi$) implique donc que (4.4.9.2) est un iso. \deleted{\unsure{Il}} Moyennant la vérification en suspens (l'existence des sommes de cardinal $\leq\pi$ dans $\operatorname{Ind}'_{\pi}(C)$), le lemme 4.4.9. résulte du résultat plus général

% lines 8-16
\deleted{\uline{Corollaire}} \inserted{Lemme} 4.4.9.0. Considérons un \deleted{diagramme} \inserted{carré commutatif} (4.4.9.1) (traits pleins) de catégories, où $\varepsilon$ admet un adjoint à gauche $\lambda$, $\varepsilon'$ un adjoint à gauche $\lambda'$, et la flèche de chgt. de base (4.4.9.2) est iso. Alors pour tout $y'_0$ dans $B'_0$, posant $\varepsilon'(y'_0)=z'_0$, le foncteur induit par $(\varepsilon,\varepsilon')$:
\[
\overline B=B_{/y'_0}\xrightarrow{\overline\varepsilon}\overline C=C_{/z'_0}
\]
admet un adjoint à gauche $\overline\lambda$, donné par

% lower composite with triangle
\diagram{\begin{tikzcd}[ampersand replacement=\&,column sep=large]
{\overline\lambda\bigl(\gamma(\deleted{\unsure{y}}z)\xrightarrow{u}z'_0\bigr)=\bigl(\lambda'\gamma(z)\simeq\beta(\underline{\lambda(z)})}
\arrow[r,"\lambda'(u)"] \arrow[ddr,"\substack{\overline\lambda(u)=\text{composé}\\\mathrm{adj}\circ\lambda'(u)}"']
\& \lambda'(z'_0) \arrow[d,equal] \\
\& \lambda'\varepsilon'(y'_0) \arrow[d,"\mathrm{adj}."] \\
\& {y'_0\bigr)}
\end{tikzcd}}
